\section{Finite-Variance Equivalence and Global Rigidity}
\label{sec:finite-variance-equivalence-and-global-rigidity}

Two losses of information drive this section.  First, the product sample space
is \(\R^{\N}\), not \(\ell^2\): row sums defined only up to separate null sets
cannot simply be composed.  We construct a common measurable carrier before
using inverse maps, tail symmetries, or likelihood cocycles.  Second, convergence
of individual remote rows gives no summability.  The necessity proof combines
qualitative row rigidity with carefully chosen tail-coordinate swaps that turn
any nonsummable error into a locally visible but asymptotically invisible
Hellinger perturbation.
Hilbert--Schmidt sufficiency, the orbit dichotomy, and the finite-marginal
criterion complete the classification.  The classification then supplies the
direct-sum compactness argument needed for global finite-dimensional
coercivity.

\input{parts/measurable-realization-and-measure-class-dichotomy}
\input{parts/finite-variance-classification}
